\section{Introduction}\label{sec:introduction}

A quadratic objective on a box depends on the first and second moments of
its variables. Replacing the possible moment vectors by their convex hull
therefore gives an exact convex formulation. The difficulty is describing
that hull. Even a relaxation built from many valid inequalities can omit a
nonnegative quadratic with only three variables. Conversely, the existence
of an exact formulation in three variables does not imply that a larger
finite semidefinite formulation will be exact in four variables.

This paper studies these two issues through the cone of quadratics
nonnegative on $[0,1]^n$. We first identify a family of valid three-variable
inequalities that recent semidefinite and second-order-cone relaxations
miss. We then give an exact, small semidefinite formulation of that family.
The same cone provides the setting for a dimension threshold: the full
quadratic moment hull has a finite semidefinite lift precisely when
$n\le3$. The negative result also holds when positive surplus is allowed
on the diagonal moments. This variant is the appropriate joint hull for
objectives whose square coefficients are nonnegative.

The connection between the two parts is concrete. A nonnegative quadratic
is a valid linear inequality for the moment hull. A certificate system
describes a subcone of these quadratics, and its dual defines a relaxation.
The new family fills a proved gap in one such system. The representability
result determines where any finite semidefinite system, regardless of its
chosen certificates or auxiliary variables, must cease to be exact.

\subsection{Relation to prior work}

Standard box relaxations combine the Shor moment-matrix condition
\citep{Shor1987QuadraticOptimization} with McCormick inequalities for
bilinear products \citep{McCormick1976}. The continuous polynomial form
of the reformulation--linearization technique (RLT) provides a systematic
framework for linearizing products of bound inequalities
\citep{SheraliTuncbilek1992}.

Anstreicher and Burer gave exact completely positive descriptions of
quadratic moment hulls on low-dimensional polytopes
\citep{AnstreicherBurer2010}. Triangulating a polytope of dimension at most
three produces completely positive blocks of order at most four. At these
orders, complete positivity is equivalent to positive semidefiniteness
and entrywise nonnegativity \citep{Diananda1962,MaxfieldMinc1962}.
Thus the exact three-variable box hull already has a finite semidefinite
lift. We use that construction both as a mathematical reference and as an
exact local strengthening in the archived computational comparisons. We do
not propose a first exact formulation for this hull.

Burer and Dong developed recursive descriptions and separation methods
for the related box completely positive and copositive cones
\citep{BurerDong2013}. An exact separation procedure and a finite
semidefinite lift are different properties. Our impossibility theorem
concerns the latter; it does not prevent separation or other exact
optimization methods.

Recent work seeks formulations that exploit sparse objectives and small
local blocks. Khajavirad's disjoint-support construction uses sums of
affine squares weighted by products of box literals, with the affine
factor involving only variables absent from the weight
\citep{Khajavirad2026SparseBoxQP}. In three variables with positive square
coefficients, the resulting extended formulation consists of $27$
localizing matrices. Anstreicher and Puges strengthen the usual
three-variable inequalities using a trilinear auxiliary variable and
second-order-cone inequalities \citep{AnstreicherPuges2025ETRI}. We give a
single rational moment functional feasible for both systems, including all
coordinate complements and permutations of the latter inequalities, but
violating a nonnegative quadratic. This is a comparison against the
specified systems, rather than a claim about all possible SOS or SOC
relaxations.

Sign structure supplies an additional point of comparison.
Burer, Natarajan, and Willemsen prove exactness of a semidefinite
relaxation for submodular box quadratics in at most three variables
\citep{BurerNatarajanWillemsen2025}. Combined with cube symmetries, this
result settles the three-variable disjoint-support certificate question
for quadratics with nonnegative square coefficients whenever the product
of the three mixed coefficients is nonpositive.
The new inequalities occupy the remaining sign class. This restriction
helps explain why they require three variables and why two-variable
restrictions do not detect them.

Hildebrand's rank-one subtraction criterion for copositive matrices
\citep{Hildebrand2014MinimalZeros} supplies a further geometric tool.
Applied through the cube vertices, it implies that the zeros of a
non-square extreme quadratic must affinely span the box. We combine this
criterion with an explicit analysis of convex facet restrictions and a
classification of cube-edge contact graphs. Together with the submodular
duality argument, these steps classify every extreme ray with positive
square coefficients and no vertex zero. The subtraction criterion is
prior work; the facet analysis and the resulting box classification are
the developments proved here.

For the negative representability result, the essential external tool is
the characterization of spectrahedral shadows by preservation under
completely positive maps over real closed fields due to Bodirsky, Kummer,
and Thom \citep{BodirskyKummerThom2024}. Their construction uses the Horn
copositive matrix to exhibit an obstruction to semidefinite
representability. We transfer that obstruction to quadratics on a box by
using the constant monomial as one copositive coordinate and placing the
other four coordinates near a box vertex. The transfer rules out every
finite semidefinite lift, rather than only a doubly nonnegative
approximation. We also extend the argument to polytopes with a simple
vertex and to sparse positive-loop graphs containing a $K_4$ minor.
These conclusions use the existing preservation theorem; the application
to the box and the stated sparse hulls is the result proved here.

Nishijima proves that the copositive and completely positive cones over
any symmetric cone of rank at least five are not spectrahedral shadows
\citep[Theorem~3.2]{Nishijima2026}. The present paper uses the explicit
box embedding and graph reductions to reach the moment hulls studied here.

\subsection{Contributions and scope}

The paper makes the following contributions.

\begin{enumerate}[leftmargin=*,itemsep=5pt]
\item An exact rational certificate proves that the three-variable
disjoint-support system misses a quadratic with strictly positive square
coefficients.
The same certificate satisfies the switched SOC strengthening of
Anstreicher and Puges. A direct proof using five edge zeros explains the
failure without a numerical optimization argument
(\cref{family:counterexample-section}).

\item A parameter family generalizes this example. Its validity follows
from an explicit polynomial identity. In an open parameter region its
members generate exposed extreme rays of the nonnegative-quadratic cone
and admit no disjoint-support certificate. All inequalities in the larger
validity domain are enforced by one positive semidefinite block of order
five and six nonnegative scalars, using only first and second moments
(\cref{family:section}).

\item We establish the closedness and dual interpretation of the
disjoint-support cone, relate diagonal caps to the distinction between
the full and positive-loop hulls, and prove the sign restriction and
boundary extreme-ray results needed to state a precise three-variable
completeness question. We further classify every extreme ray with positive
square coefficients and positive values at all cube vertices, reducing
the remaining completeness question to vertex-zero rays
(\cref{sec:three}). The completeness question remains separate
from the validity and enforcement theorems.

\item We prove the sharp dimension threshold for finite semidefinite
lifts of the full box moment hull, the positive-loop hull, and their
dual cones. We prove a local geometric extension at simple vertices and
show that every proper face of the four-variable box hull nevertheless
has such a lift (\cref{repr:section}).

\item For sparse moment hulls, we prove a $K_4$-minor obstruction on the
positive-loop subgraph and construct finite lifts when every positive-loop
component has at most three vertices. The construction gives an explicit
size bound in terms of a tree decomposition after the component boundaries
are made cliques (\cref{graph:section}). Archived computations compare the
new inequalities with the named relaxations and exact local hulls, with
separate conclusions for constructed and benchmark instances
(\cref{comp:section}).
\end{enumerate}

To the best of our knowledge, the family and its order-five formulation,
the exact separating certificate against the two named recent systems,
the classification of three-variable extreme rays with positive square
coefficients and no vertex zero, and the
stated box and sparse-graph representability consequences have
not appeared in the prior work reviewed here. The proofs below specify
these claims independently of priority. In particular, neither the
classical order-four cone identity nor the existing three-variable
tetrahedral lift is a new result.

The practical use of the family is local convexification. Each selected
triple contributes a block with no trilinear or higher-order moment
variables, so the inequalities can strengthen a relaxation that already
shares first and second moments. The archived study contains constructed
instances where this formulation attains the bound of exact triple hulls
with fewer selected blocks. On the tested natural instances, the base
relaxations were often already exact, and the new family produced little
additional improvement. These findings support selective use of the
family; they do not establish a general speed advantage or a
branch-and-bound performance improvement.

\Cref{sec:setting} fixes the moment and coefficient conventions.
\Cref{family:counterexample-section,family:section,sec:three} develop the
three-variable inequalities and their geometry.
\Cref{repr:section,graph:section} treat finite representability.
\Cref{comp:section} presents the archived computational evidence.
The appendices contain the exact moment certificate, boundary calculations,
stationary-facet and contact-graph proofs, real closed field details, and
separation and comparison tools.
